\paragraph{跨接口终端结论：视界、谱与纤维几何的可复用陈述}
本段把全文中已建立的若干终端断言，以“可审计、可独立引用”的格式汇总为结论条目；每条结论均对应正文或附录中的已证明命题/定理（在条目末给出引用），因此不再重复证明细节。

\begin{conclusion}[视界证书复杂度的二次--线性双下界：Toeplitz 阶数与共动扫描覆盖均不可避免发散]\label{con:xi-terminal-horizon-certificate-quadratic-linear-lb}
固定偏移下界 \(\delta\ge \delta_0>0\) 与高度窗 \(|\gamma|\le T\)。
若要求制度内证书在该窗内对任意可能的缺陷位置给出一致可检测性，则无论采用 Toeplitz--PSD 阶数提升（矩层级），或采用共动扫描的有限覆盖（有限读出族），其资源规模均必须随 \(T\) 发散，并满足显式下界
\[
N_{\mathrm{Toep}}\ \gtrsim\ Q_{\max}(T,\delta_0)\,\log\!\bigl(1+\delta_0 Q_{\max}(T,\delta_0)\bigr)
\asymp \frac{T^2}{\delta_0}\,\log(1+T^2),
\qquad
N_{\mathrm{scan}}\ \gtrsim\ \frac{T}{t+\delta_0}\quad(t>0\ \text{固定}),
\]
其中
\[
Q_{\max}(T,\delta_0):=\frac{T^2+(1+\delta_0)^2}{4\delta_0}.
\]
\end{conclusion}
\begin{proof}[证明]
见命题 \ref{con:xi-toeplitz-vs-scan-complexity-lower-bounds}（并结合其引用的推论/定理）。
\end{proof}

\begin{conclusion}[负谱幅度读出边界层深度：有限 Toeplitz 否证器给出缺陷贴边尺度]\label{con:xi-terminal-negative-spectrum-boundary-layer-depth}
设某阶 Toeplitz 主子矩阵 \(\mathbf{T}_N\) 的最小特征值为 \(\lambda_{\min}(\mathbf{T}_N)=-\eta_N<0\)。
则存在盘内零点 \(a\in\mathbb D\) 使其边界层深度 \(h:=1-|a|\) 满足
\[
h\ \ge\ c\,\frac{\eta_N}{N},
\]
其中 \(c>0\) 为绝对常数。
因此 \(\eta_N\) 在制度语义中不仅是“是否失败”的二值否证器，而且给出缺陷贴边深度的定量证书；其自然归一化量纲为 \(\eta_N/N\)。
\end{conclusion}
\begin{proof}[证明]
见命题 \ref{con:xi-toeplitz-negative-eigenvalue-controls-depth}。
\end{proof}

\begin{conclusion}[共动地址回收的位预算律：\texorpdfstring{$2\log_2T$}{2log2T} 主阶高度依赖可被制度性消除]\label{con:xi-terminal-comoving-address-bit-budget-recovery}
固定 \(\delta_0\in(0,\tfrac12)\) 与高度窗 \(|\gamma|\le T\)。
若不外置共动前缀，则端点层最小可分辨余量 \(2^{-p}\) 的一致排除要求强制
\[
p\ \ge\ 2\log_2 T+\log_2\!\frac{1}{\delta_0}+O(1)\qquad (T\to\infty).
\]
若允许外置可寻址前缀并取共动局部化 \(x_0=\gamma\)，则可检测余量满足高度无关的统一下界
\[
h_{\gamma}(\rho)=1-|w_{\rho,\gamma}|^2=\frac{4\delta}{(1+\delta)^2}\ \ge\ \frac{4\delta_0}{(1+\delta_0)^2},
\]
从而位预算门槛退化为常数下界
\[
p\ \ge\ \log_2\!\Bigl(\frac{(1+\delta_0)^2}{4\delta_0}\Bigr),
\]
并由两式对照读出：共动前缀在高度窗极限中制度性回收的主阶位预算为 \(2\log_2T+O(1)\)。
\end{conclusion}
\begin{proof}[证明]
见定理 \ref{thm:xi-offcritical-endpoint-resolution-lower-bound} 与推论 \ref{cor:xi-comoving-prefix-bit-budget-gain}。
\end{proof}

\begin{conclusion}[Minkowski 预算障碍：最坏扭曲能量的普适上界为 \texorpdfstring{$B^{-2/d}$}{B^{-2/d}}]\label{con:xi-terminal-minkowski-budget-barrier}
在 \(d\)-维扭曲/角色预算中，令 \(B=\prod_k m_k\) 为总模预算。
若最短向量落入小角有效域，则最坏扭曲能量 \(q(m)\) 满足普适上界
\[
q(m)\ \le\ (2\pi)^2\,C_d\,(\det\Sigma)^{1/d}\,B^{-2/d},
\qquad
C_d:=\Bigl(\frac{2^{d}}{V_d}\Bigr)^{2/d},
\]
其中 \(V_d\) 为单位球体积。
因此在固定维数 \(d\) 下，任何试图以增大模预算改善最坏扭曲能量的策略，至多获得指数 \(B^{-2/d}\) 的幂律改进；该指数为结构性上界而非算法松弛。
\end{conclusion}
\begin{proof}[证明]
见定理 \ref{thm:pom-minkowski-budget-barrier}。
\end{proof}

\begin{conclusion}[回溯熵缺口的指数放大：ordinary 与 non-backtracking 的增长率差异为 \texorpdfstring{$e^{n\Delta h}$}{exp(n Dh)}]\label{con:xi-terminal-backtracking-entropy-gap-exponential}
对无向骨架的邻接谱半径 \(\rho(A)\) 与 Hashimoto 无回溯谱半径 \(\rho(B)\)，定义回溯熵缺口
\[
\Delta h:=\log\rho(A)-\log\rho(B)=\log\!\Bigl(\frac{\rho(A)}{\rho(B)}\Bigr)\ge 0.
\]
则 ordinary closed walk 与 non-backtracking closed walk 的增长率比在长度 \(n\) 上满足
\[
\Bigl(\frac{\rho(A)}{\rho(B)}\Bigr)^n=e^{n\Delta h}.
\]
在单流实例中，表 \ref{tab:parallel_addition_kernels_backtracking_entropy_gap} 给出 \(n=50\) 的放大量级（其中 \(21\)-local 的放大已达 \(10^2\) 量级），从而 Ihara primes（无回溯）正规化在量化层面成为必要步骤。
\end{conclusion}
\begin{proof}[证明]
见注 \ref{rem:ihara-entropy-gap} 与表 \ref{tab:parallel_addition_kernels_backtracking_entropy_gap} 的定义与审计输出。
\end{proof}

\begin{conclusion}[模障碍当且仅当 cyclotomic 因子：Ihara--Artin 判别在有限维代数层面闭合]\label{con:xi-terminal-mod-obstruction-cyclotomic}
固定 \(u>0\)。
对每个不可约表示 \(\rho\neq\mathbf 1\) 的扭曲 Hashimoto 算子 \(B_\rho(u)\)，令
\[
Z_K(u,t;\rho):=\det(I-tB_\rho(u))^{-1}.
\]
则存在非平凡模障碍当且仅当存在某个 \(\rho\neq\mathbf 1\) 使 \(\rho(B_\rho(u))=\rho(B_{\mathbf 1}(u))=:\lambda_1(u)\)；并且在该情形下 \(\det(I-tB_\rho(u))\) 必含 cyclotomic 因子（等价地，归一化多项式 \(\det(I-(t/\lambda_1(u))B_\rho(u))\) 在单位圆上出现单位根因子）。
反之，若某 \(\det(I-tB_\rho(u))\) 出现 cyclotomic 因子，则必有 \(\rho(B_\rho(u))=\lambda_1(u)\)。
在广义 Schur/Carath\'eodory 的 Pontryagin 实现口径下（\S\ref{subsubsec:xi-negative-square-index-pontryagin-realization}），
上述 cyclotomic 因子亦可等价表述为：归一化算子 \(\lambda_1(u)^{-1}B_\rho(u)\) 在单位圆上出现单位根特征值，
且其特征子空间在实现所诱导的不定型内积上为中性（各向同性，零签名）。
因此模障碍并非“频率偶然命中”，而是单位根处中性点谱的结构性碰撞；其可枚举性来自 cyclotomic 因子的完全分类。
\end{conclusion}
\begin{proof}[证明]
见定理 \ref{thm:conclusion75-ihara-cyclotomic-factor-classification}。
\end{proof}

\begin{conclusion}[Big-Witt/necklace 桥包的严格交换与函子性：\texorpdfstring{$\zeta=\mathrm{Euler}\circ\mathrm{Mob}\circ\mathrm{TrSeq}$}{zeta=Euler o Mob o TrSeq}]\label{con:xi-terminal-big-witt-functoriality}
对任意有限核/有限矩阵 \(T\)，其迹序列（ghost 坐标）
\(
a_n(T):=\Tr(T^n)
\)
在 intertwiner/相似变换下保持不变；Möbius/Witt 反演给出 primitive 坐标 \(p_n(T)=\frac1n\sum_{d\mid n}\mu(d)\,a_{n/d}(T)\)。
并且成立逐项恒等式
\[
\zeta_T(z):=\det(I-zT)^{-1}
=\exp\!\Bigl(\sum_{n\ge 1}\frac{a_n(T)}{n}z^n\Bigr)
=\prod_{n\ge 1}(1-z^n)^{-p_n(T)}.
\]
因此 primitive 频谱（Witt 坐标）仅依赖于 \(T\) 的 intertwiner 类而与具体呈示无关。
\leanverified{paper\_zeta\_necklace\_counting\_seeds}
\end{conclusion}
\begin{proof}[证明]
见注 \ref{rem:zetaK-commuting-diagram}；带 Adams 操作与参数函子化版本见 \S\ref{subsec:zeta-online-kernel}。
\end{proof}

\begin{conclusion}[window-\(6\) 边界态的 \texorpdfstring{$\ZZ/2$}{Z/2} sheet parity：二层别名门户由 \texorpdfstring{$F_9=34$}{F9=34} 唯一承载]\label{con:xi-terminal-window6-bdry-sheet-parity}
令 window-\(6\) 的 boundary 扇区
\(
X^{\mathrm{bdry}}_6:=\{w\in X_6:\ w_1=w_6=1\}
\)。
则对任意 \(w\in X^{\mathrm{bdry}}_6\) 有严格二层别名
\[
(\mathrm{Fold}^{\mathrm{bin}}_6)^{-1}(w)=\{V_6(w),\,V_6(w)+34\},
\]
从而 dyadic 预像在组合层面被限制为一对相距 \(F_9=34\) 的双层，实现内生的 \(\ZZ/2\) sheet parity。
\end{conclusion}
\begin{proof}[证明]
见推论 \ref{cor:terminal-foldbin6-boundary-pure-f9-alias}。
\end{proof}

\begin{conclusion}[黄金互补二元信道的最强数据处理常数：\texorpdfstring{$\delta=\rho_{\mathrm{m}}=\varphi^{-3}$}{delta=rho_m=phi^-3} 与 \texorpdfstring{$I$}{I}-收缩 \texorpdfstring{$\varphi^{-6}$}{phi^-6}]\label{con:xi-terminal-golden-binary-channel-sdpi}
取二元信道 \(W\) 使 \(p=\varphi^{-1}\)、\(q=\varphi^{-2}=1-p\)，并记 \(\Delta=p-q=\varphi^{-3}\)。
则
\[
\delta(W)=\mathrm{TV}(\mathrm{Bern}(p),\mathrm{Bern}(q))=|\Delta|=\varphi^{-3},
\qquad
\rho_{\mathrm{m}}(W)=|\Delta|=\varphi^{-3}.
\]
从而对任意马尔可夫链 \(U\to \Theta\to B\) 成立锐常数的最强数据处理不等式
\[
I(U;B)\ \le\ \rho_{\mathrm{m}}(W)^2\,I(U;\Theta)\ =\ \varphi^{-6}\,I(U;\Theta),
\]
且 \(\varphi^{-6}\) 为不可改进的收缩因子。
\end{conclusion}
\begin{proof}[证明]
见定理 \ref{thm:pom-golden-binary-channel-contraction}。
\end{proof}

\begin{conclusion}[位串 \texorpdfstring{$\to$}{->} 有理证书的可回放性：半角连分数正规化给出唯一证书或 \texorpdfstring{$\mathsf{NULL}$}{NULL}]\label{con:xi-terminal-bitstring-rational-certificate-replayability}
给定浮点位串输入 \(\widehat u\in\CC\)、最大分母 \(Q_{\max}\in\NN\) 与误差预算 \(\delta>0\)，
通过舍入前像区间 \(\mathcal I(\widehat u)\)、半角变量与确定性连分数截断，可输出唯一证书 \((p,q)\)（从而 \(u_{\mathrm{cert}}=r(p,q)\in\TT\)）或空值 \(\mathsf{NULL}\)。
并且接受判据 \(\mathrm{ACCEPT}(p,q)\) 仅依赖 \((p,q)\)、\(\mathcal I(\widehat u)\)、\(\delta\) 与固定范数上界计算；第三方仅凭位串与输出即可复核真值，从而满足严格可回放性与跨平台一致性。
\end{conclusion}
\begin{proof}[证明]
见算法 \ref{alg:unitary-two-integer-deterministic} 与命题 \ref{prop:unitary-two-integer-reproducible}。
\end{proof}

\begin{conclusion}[素长 primitive 的 Wieferich 指纹：\texorpdfstring{$p$}{p}-进退化强度由单一整数 \texorpdfstring{$\nu_p$}{nu_p} 完全分级]\label{con:xi-terminal-wieferich-fingerprint}
令 \(p\) 为奇素数，\(L_p\) 为 Lucas 数，并定义素长 Wieferich 指纹
\(
\nu_p:=v_p(L_p-1)\in\ZZ_{\ge 1}
\)。
则真实输入 40 状态核的素长 primitive 轨道数 \(p_p\) 满足
\[
v_p(p_p)=\nu_p-1,
\]
从而对任意 \(k\ge 1\) 有整数判据
\[
p^k\mid p_p\ \Longleftrightarrow\ \nu_p\ge k+1.
\]
因此素长层的“算术退化强度”由单一整数 \(\nu_p\) 完全分级：\(\nu_p=1\) 为非退化相，而 \(\nu_p\ge 2\) 对应 Lucas--Wieferich 型稀有退化相；并且退化阶数与 \(p\)-进赋值精确对齐，从而 \(\nu_p\) 可作为该退化机制的单参序参量。
特别地，\(p\mid p_p\) 当且仅当 \(p^2\mid(L_p-1)\)，即素长层的 primitive 算术退化精确等价于 Lucas--Wieferich 型事件。
\end{conclusion}
\begin{proof}[证明]
见定义 \ref{def:pom-lucas-nu-p} 与推论 \ref{cor:pom-lucas-nu-p-criterion}。
\end{proof}

\begin{conclusion}[primitive 奇偶分裂的结构性消项：\texorpdfstring{$\lambda^{n/2}$}{lambda^{n/2}} 次项在偶长上精确抵消，首个可见振荡降至 \texorpdfstring{$\lambda^{n/4}$}{lambda^{n/4}}]\label{con:xi-terminal-primitive-parity-cancellation}
记 \(\lambda=\varphi^2\) 为主增长率，\(p_n\) 为 primitive 轨道数。
则存在奇偶分裂的渐近展开
\[
p_n=\frac{\lambda^n}{n}+\frac{\bigl((-1)^n-\mathbf{1}_{2\mid n}\bigr)\lambda^{n/2}}{n}
+O\!\left(\frac{\lambda^{n/3}}{n}\right),
\]
其中偶数情形的 \(\lambda^{n/2}\) 级项在 Möbius 反演的 \(d=1,2\) 两项间发生精确抵消；因此当 \(n\) 为偶数时该次项制度性消失。
进一步在偶长 \(n=2m\) 情形，首个可见振荡次项来自达界特征值 \(-\sqrt{\lambda}\)，并满足
\[
p_{2m}=\frac{\lambda^{2m}}{2m}-\frac{(-\sqrt{\lambda})^{m}}{2m}+O\!\left(\frac{\lambda^{2m/3}}{m}\right),
\]
其规模为 \(\lambda^{n/4}\)。
\end{conclusion}
\begin{proof}[证明]
见推论 \ref{cor:pom-prim-odd-even} 与注 \ref{rem:pom-prim-even-lower-osc}。
\end{proof}

\begin{conclusion}[Fibonacci 张量闭包类的 RH-sharp 生成器：平方根界在 \texorpdfstring{$(-\varphi)^n$}{(-phi)^n} 上制度性饱和]\label{con:xi-terminal-fib-tensor-rhsharp-generator}
令
\(
F=\bigl(\begin{smallmatrix}1&1\\[2pt]1&0\end{smallmatrix}\bigr)
\)、\(B:=F\otimes F\)、\(C:=-F\)。
对任意 \(n\ge 1\) 与非负整数 \(a,b\)，定义
\[
M_{n;a,b}:=B^{\otimes n}\ \oplus\ C^{\otimes n}\ \oplus\ I_a\ \oplus\ (-I)_b,
\qquad
\lambda_n:=\rho(M_{n;a,b}).
\]
则 \(\lambda_n=\varphi^{2n}\)，并且对任意非主特征值 \(\mu\in\sigma(M_{n;a,b})\) 有平方根界 \(|\mu|\le \sqrt{\lambda_n}=\varphi^n\)，且该界在 \(C^{\otimes n}\) 的主特征值 \(\mu=(-\varphi)^n\) 上取等号。
此外，\(\Tr(M_{n;a,b}^k)\) 与 primitive 轨道数具有 Lucas--Möbius 的显式闭式表达。
\end{conclusion}
\begin{proof}[证明]
见定理 \ref{thm:fib-tensor-rhsharp-closure}。
\end{proof}

\begin{conclusion}[中心化三阶迹的基线剥离与 OPUC/CMV 参数化：大值 Gram 的反射谱化]\label{con:xi-terminal-centered-trace3-opuc-cmv}
令 \(G\) 为大值点集诱导的 Gram 矩阵，并令
\(
G^\circ:=G-\frac{\tr(G)}{|W|}I
\)。
则中心化三阶迹满足恒等式
\[
\tr\!\bigl((G^\circ)^3\bigr)
=\tr(G^3)-3\frac{\tr(G)}{|W|}\tr(G^2)+2\frac{\tr(G)^3}{|W|^2},
\]
并且在 \(R(v)=\sum_{t\in W}v^{it}\) 的表示下，与 \(\int |R|^4\) 同阶的四体项可被严格替换为方差型能量 \(E(W)-c|W|^2\)，从而中心化机制剥离确定性随机基线，仅对超随机能量部分敏感。
进一步，在带宽窗口内，\(G\) 可视为离散测度的有限截断 Toeplitz 抽样压缩；由 OPUC 理论存在唯一 Verblunsky 系数序列与 CMV 矩阵，使中心化迹量可表示为 Verblunsky 系数的有限多项式泛函并受有限节谱测度控制。
\end{conclusion}
\begin{proof}[证明]
见命题 \ref{prop:gm-centered-trace3-energy-variance} 与命题 \ref{prop:gm-toeplitz-cmv-verblunsky-embedding}。
\end{proof}

\begin{conclusion}[Sofic 扭曲算子的张量闭合：统一谱隙 \texorpdfstring{$\Rightarrow$}{=>} 全奇阶中心化迹指数压缩 \texorpdfstring{$\Rightarrow$}{=>} 可调阶矩把大值频次压回均方]\label{con:xi-terminal-sofic-tensor-twist-moment-compression}
对 sofic--Zeckendorf 转导器生成的读出模型，中心化迹量 \(\tr((G^\circ)^\ell)\) 可表达为有限多个张量扭曲转移算子谱控制的线性组合。
若对所有非平凡扭曲存在统一谱隙 \(\rho(A(\chi))\le (1-\kappa)\rho(A(1))\)，则存在 \(c_\ell(\kappa)>0\) 使
\[
|\tr((G^\circ)^\ell)|\ \ll\ \rho(A(1))^{\ell m}\,e^{-c_\ell(\kappa)m},
\]
并由此得到可调阶矩极限：在统一谱隙下可选取阶数 \(k=k(\varepsilon,\kappa)\) 使归一化矩阵谱范数满足 \(\|\widetilde G\|_{\mathrm{op}}\le 1+O(|W|^{-\varepsilon})\)，从而关键区间的大值频次被压回均方并获得严格小幂次修正。
\end{conclusion}
\begin{proof}[证明]
见定理 \ref{thm:gm-tensor-twist-centered-moments} 与推论 \ref{cor:gm-adjustable-moment-mean-square}。
\end{proof}

\begin{conclusion}[稳定层因子链的谱--熵双下界：混合、超收缩与次高斯浓缩的统一常数化]\label{con:xi-terminal-factor-chain-gap-lsi-hypercontractive}
在折叠因子链 \(P^{\mathrm{fac}}_m\)（平稳分布 \(\pi_m\)）上，谱隙与对数 Sobolev 常数满足统一下界
\[
\mathrm{gap}(P^{\mathrm{fac}}_m)\ge \frac{2}{m},
\qquad
\alpha(P^{\mathrm{fac}}_m)\ge \frac{1}{m}.
\]
因此导通率下界、Gross 超收缩以及在平方跳变能量归一化下的次高斯浓缩均可被下界化为超立方体尺度 \(m^{-1}\)，且与纤维多重度的局部粗糙性无关。
\end{conclusion}
\begin{proof}[证明]
见定理 \ref{thm:pom-fold-factor-chain-gap-lsi} 与推论 \ref{cor:pom-fold-factor-chain-derived-invariants}。
\end{proof}

\begin{conclusion}[无自环刚性：一步超立方体翻转必改变稳定类型，故因子链为严格零对角核]\label{con:xi-terminal-factor-chain-no-self-loop}
折叠因子链满足
\[
P^{\mathrm{fac}}_m(x,x)=0,\qquad \forall x\in X_m,
\]
其结构来源为：超立方体单比特翻转在微态空间上诱导的差分与同一稳定类型纤维内允许的差分不相容，从而不可能发生一步“停留”。
\end{conclusion}
\begin{proof}[证明]
见命题 \ref{prop:pom-fold-factor-chain-reversible-noloop}。
\end{proof}

\begin{conclusion}[纤维重构图的 Cartesian 立方闭包：维数与直径的完全闭式]\label{con:xi-terminal-fiber-reconstruction-cartesian}
对任意 \(x\in X_m\)，纤维重构图 \(\mathcal F_x\) 在 \(\Gamma_x\simeq\bigsqcup_i P_{\ell_i}\) 的分量长度 \((\ell_i)\) 上严格分解为 Fibonacci 立方的 Cartesian 乘积：
\[
\mathcal F_x\ \cong\ \prod_i \Gamma_{\ell_i}.
\]
并且其最大立方维数与直径具有闭式
\[
\dim_{\square}(\mathcal F_x)=\sum_i \left\lceil\frac{\ell_i}{2}\right\rceil,
\qquad
\mathrm{diam}(\mathcal F_x)=\sum_i \ell_i .
\]
\end{conclusion}
\begin{proof}[证明]
见定理 \ref{thm:pom-fiber-reconstruction-cartesian-product}。
\end{proof}
\begin{theorem}[Fold 纤维规模的 \texorpdfstring{$\varphi$}{phi}-指数律]\label{thm:xi-fold-fiber-fibonacci-phi-exponential-law}
\leanverified{paper\_xi\_fold\_fiber\_fibonacci\_phi\_exponential\_law}
令 \(\varphi=\frac{1+\sqrt5}{2}\)。
对任意 \(m\ge 1\) 与任意 \(x\in X_m\)，令候选位点诱导图分解为
\(
\Gamma_x\simeq\bigsqcup_{i=1}^{r}P_{\ell_i}
\)。
记纤维基数 \(d_m(x)=\card{\Fold_m^{-1}(x)}\)，并沿用结论 \ref{con:xi-terminal-fiber-reconstruction-cartesian} 的直径与立方维数记号。
则存在绝对常数使得以下两组不等式同时成立：
\[
\varphi^{\mathrm{diam}(\mathcal F_x)}\ \le\ d_m(x)\ \le\ \varphi^{\mathrm{diam}(\mathcal F_x)+r},
\]
\[
\varphi^{2\dim_{\square}(\mathcal F_x)-r}\ \le\ d_m(x)\ \le\ \varphi^{2\dim_{\square}(\mathcal F_x)+r}.
\]
等价地，
\[
\Bigl|\log d_m(x)-2(\log\varphi)\,\dim_{\square}(\mathcal F_x)\Bigr|\ \le\ r\log\varphi.
\]
\end{theorem}
\begin{proof}
由定理 \ref{thm:pom-fiber-indset-factorization}，有乘积表达式
\(
d_m(x)=\prod_{i=1}^{r}F_{\ell_i+2}
\)，其中 \(F_k\) 为 Fibonacci 数。
标准估计给出对一切 \(k\ge 2\) 有 \(\varphi^{k-2}\le F_k\le \varphi^{k-1}\)，故
\[
\varphi^{\ell_i}\ \le\ F_{\ell_i+2}\ \le\ \varphi^{\ell_i+1}.
\]
逐项相乘得
\[
\varphi^{\sum_i \ell_i}\ \le\ d_m(x)\ \le\ \varphi^{\sum_i \ell_i+r}.
\]
再由结论 \ref{con:xi-terminal-fiber-reconstruction-cartesian} 的恒等式 \(\mathrm{diam}(\mathcal F_x)=\sum_i\ell_i\) 得第一组不等式。
\par
对第二组不等式，仅需注意对任意整数 \(\ell\ge 1\) 有
\(
2\lceil \ell/2\rceil-1\le \ell\le 2\lceil \ell/2\rceil
\)。
对 \(\ell_i\) 求和并用 \(\dim_{\square}(\mathcal F_x)=\sum_i\lceil \ell_i/2\rceil\) 得
\[
2\dim_{\square}(\mathcal F_x)-r\ \le\ \mathrm{diam}(\mathcal F_x)\ \le\ 2\dim_{\square}(\mathcal F_x).
\]
将其代入第一组不等式即得第二组不等式与对数形式。证毕。
\end{proof}
\begin{conclusion}[纤维几何的完备同构谱：分量长度多重集为完全不变量，并决定对称群阶]\label{con:xi-terminal-fiber-geometry-complete-spectrum}
在 \(\mathcal F_x\cong\prod_i \Gamma_{\ell_i}\) 的设定下，分量长度多重集
\[
\mathsf{Spec}_{\square}(x):=\{\ell_1,\dots,\ell_r\}
\]
是 \(\mathcal F_x\) 的完备图同构不变量：若 \(\mathcal F_x\cong \mathcal F_y\) 则 \(\mathsf{Spec}_{\square}(x)=\mathsf{Spec}_{\square}(y)\)（多重集意义下）。
并且全部对称性仅来源于逐因子 \(\ZZ_2\) 反演与同长度因子的置换，从而 \(|\mathrm{Aut}(\mathcal F_x)|\) 由 \(\mathsf{Spec}_{\square}(x)\) 完全决定。
\end{conclusion}
\begin{proof}[证明]
见定理 \ref{thm:pom-fibcube-aut-z2}、定理 \ref{thm:pom-fiber-reconstruction-aut-group} 与推论 \ref{cor:pom-fiber-reconstruction-spectrum-complete}。
\end{proof}

\begin{conclusion}[Gray--Hamilton 结构的纤维化：全预像存在一位开关的常延迟遍历]\label{con:xi-terminal-gray-hamilton-fiberization}
Fibonacci 立方 \(\Gamma_n\) 存在显式 Gray--Hamilton Hamilton 路径，覆盖全部顶点且相邻仅差一位。
结合纤维重构图的 Cartesian 分解，对任意 \(x\) 的预像纤维（等价于 \(\mathcal F_x\) 的顶点集）存在一条遍历全体预像的路径，且相邻两预像仅作一次合法开关（对称差大小为 \(1\)），从而全预像枚举可在固定局部操作代价下实现逐步回放。
\end{conclusion}
\begin{proof}[证明]
见定理 \ref{thm:pom-fibcube-gray-hamilton} 与推论 \ref{cor:pom-fiber-constant-delay-traversal}。
\end{proof}

\begin{conclusion}[纤维独立集复形的同伦二分律：要么可缩，要么为唯一维度的同伦球面]\label{con:xi-terminal-fiber-indcomplex-homotopy-dichotomy}
令 \(K_x\) 为候选位点诱导图的独立集复形。
则其几何实现满足严格二分律：若某分量长度 \(\ell_i\equiv 1\ (\mathrm{mod}\ 3)\)，则 \(|K_x|\) 可缩；否则 \(|K_x|\) 同伦等价于单一球面 \(S^{\tau(x)-1}\)，且约化同调呈“唯一维度、秩为 \(1\)”的刚性。
\end{conclusion}
\begin{proof}[证明]
见定理 \ref{thm:pom-fiber-independence-complex-classification} 与推论 \ref{cor:pom-fiber-independence-complex-homology}。
\end{proof}

在纤维独立集复形的球面--可缩二分律之外，全部持久拓扑量事实上都被进一步压缩为唯一维度上的秩一对象。

\begin{theorem}[纤维拓扑的秩一律]\label{thm:xi-time-part10-fiber-topology-rankone-law}
\leanverified{paper\_xi\_time\_part10\_fiber\_topology\_rankone\_law}
沿用定理 \ref{thm:pom-fiber-independence-complex-classification}、
推论 \ref{cor:pom-fiber-independence-complex-homology} 与定理
\ref{thm:killo-fold-fiber-homotopy-parity-equivalence} 的记号。对任意稳定类型 \(x\in X_m\)，记
\[
K_x:=\mathrm{Ind}^{\Delta}(\Gamma_x).
\]
则成立：
\begin{enumerate}
  \item 若 \(|K_x|\) 非可缩，则
  \[
  |K_x|\simeq S^{\tau(x)-1}.
  \]
  \item 因而其约化同调满足
  \[
  \widetilde H_j\bigl(|K_x|;\mathbb Z\bigr)=0
  \qquad (j\neq \tau(x)-1),
  \]
  \[
  \widetilde H_{\tau(x)-1}\bigl(|K_x|;\mathbb Z\bigr)\cong \mathbb Z.
  \]
  \item 并且
  \[
  |K_x|\not\simeq *
  \iff
  d_m(x)\ \text{为奇数}.
  \]
\end{enumerate}
特别地，每条持久纤维拓扑障碍都只能在唯一维度上以秩 \(1\) 出现；不存在同一纤维同时承载两条线性独立的持久拓扑通道。
\end{theorem}
\begin{proof}
第 \(1\) 条正是定理 \ref{thm:pom-fiber-independence-complex-classification}
在非可缩情形下的结论。第 \(2\) 条则由推论
\ref{cor:pom-fiber-independence-complex-homology} 立即得到。第 \(3\) 条正是定理
\ref{thm:killo-fold-fiber-homotopy-parity-equivalence} 的等价判别。

因此，一旦 \(|K_x|\) 非可缩，其约化同调只能在唯一维度
\(\tau(x)-1\) 上非零，且该同调群秩恰为 \(1\)。这就排除了同一纤维上同时出现两条线性独立持久拓扑通道的可能性。证毕。
\end{proof}

\endinput


\begin{conclusion}[折叠条件期望的 Rényi 平坦性：最坏信息损失与 Rényi 阶数无关且由最大纤维实现]\label{con:xi-terminal-renyi-flatness-logM}
设 \(\Omega\twoheadrightarrow X\) 为有限交换折叠，\(E:\CC^\Omega\to \CC^X\) 为按纤维均匀平均的条件期望，
记最大纤维多重度 \(M:=\max_{x\in X}|\Fold^{-1}(x)|\)。
则对任意 \(\alpha\in[1,\infty]\)，交换情形下的 sandwiched Rényi 相对熵满足平坦极值律
\[
\sup_{\rho}\,D_\alpha(\rho\|\rho\circ E)=\log M,
\]
且上确界由最大纤维上一点态同时实现（对全部 \(\alpha\) 同时取等）。
\end{conclusion}
\begin{proof}[证明]
由推论 \ref{cor:op-algebra-renyi-flatness-sup-equals-log-index} 得 \(\sup_\rho D_\alpha(\rho\|\rho\circ E)=\log\mathrm{Ind}(E)\)；
在交换折叠口径下 \(\mathrm{Ind}(E)=M\)（见命题 \ref{prop:fold-condexp-index-maxfiber}），并由交换情形的极值态刻画得到可达性（推论 \ref{cor:op-algebra-commutative-dmax-extremizers-maxfiber}）。
\end{proof}

\begin{conclusion}[查询--外置信息联合口径下的全恢复强逆：查询比例 \texorpdfstring{$\beta$}{beta} 与外置信息密度 \texorpdfstring{$b$}{b} 的最优成功指数为截断线性律]\label{con:xi-terminal-linear-strong-converse-query-sideinfo}
设 \(X\) 为有限集合（\(|X|=n\ge 2\)），\(|\Omega|=N\)，给定纤维谱 \(d:X\to\NN\) 满足 \(\sum_{x\in X}d(x)=N\)。
令
\[
\mathfrak{F}(d):=\left\{F:\Omega\to X:\ |F^{-1}(x)|=d(x)\ \ \forall x\in X\right\},
\qquad
F\sim\mathrm{Unif}(\mathfrak{F}(d)),
\]
并记 \(w(x):=d(x)/N\)、\(H(w):=-\sum_{x\in X}w(x)\log w(x)\)。
进行 \(k=\lfloor \beta N\rfloor\) 次互异查询（允许自适应），得到轨迹 \(T_{1:k}\)。
随后允许额外获得 \(B=\lfloor bN\rfloor\) 比特外置信息 \(\mathrm{enc}(F)\in\{0,1\}^B\)，并以 \((T_{1:k},\mathrm{enc}(F))\) 为输入尝试精确恢复整张映射 \(F\)。
令最优成功率为 \(\mathrm{Succ}^{\mathrm{opt}}_{\beta,b}(N)\)。
若 \(n\log N=o(N)\)，则极限存在并满足
\[
\boxed{\;
\lim_{N\to\infty}\frac{1}{N}\log \mathrm{Succ}^{\mathrm{opt}}_{\beta,b}(N)
=
\min\bigl\{0,\ b\log 2-(1-\beta)H(w)\bigr\}.\;}
\]
从而临界密度为
\[
b_{\mathrm{crit}}(\beta):=\frac{(1-\beta)H(w)}{\log 2}.
\]
当 \(b<b_{\mathrm{crit}}(\beta)\) 时成功率以精确指数率 \((1-\beta)H(w)-b\log 2\) 衰减；当 \(b\ge b_{\mathrm{crit}}(\beta)\) 时成功率不呈指数衰减（指数率为 \(0\)）。
\end{conclusion}
\begin{proof}[证明]
对给定轨迹 \(T_{1:k}\)，记 \(C_k(x)\) 为查询到的值 \(x\) 的出现次数，并令
\[
\mathcal{M}_k
:=
\frac{(N-k)!}{\prod_{x\in X}(d(x)-C_k(x))!}
\]
为在该轨迹下与 \(d\) 相容的后验可行映射数（剩余 \(N-k\) 个点的多项式计数）。
由于 \(F\) 在 \(\mathfrak{F}(d)\) 上均匀，条件于 \(T_{1:k}\) 的后验在该集合上仍为均匀。
对任意长度 \(B\) 的确定性编码 \(\mathrm{enc}(F)\)，其在后验可行集合上诱导一个至多 \(2^B\) 个码字的分割；于是最优条件成功率为
\[
\mathrm{Succ}^{\mathrm{opt}}_{\beta,b}(N\mid T_{1:k})
=\min\Bigl(1,\frac{2^B}{\mathcal{M}_k}\Bigr),
\]
从而
\(
\mathrm{Succ}^{\mathrm{opt}}_{\beta,b}(N)=\mathbb{E}\bigl[\min(1,2^B/\mathcal{M}_k)\bigr]
\)。
\par
由 Stirling 展开与 \(n\log N=o(N)\)，并记剩余经验分布
\[
R_N(x):=\frac{d(x)-C_k(x)}{N-k}\in\Delta(X),
\]
则有
\(
\frac{1}{N}\log\mathcal{M}_k=(1-\beta)H(R_N)+o_{\mathbb{P}}(1)
\)。
又由超几何大数律 \(R_N\to w\)（概率收敛），得
\(
\frac{1}{N}\log\mathcal{M}_k=(1-\beta)H(w)+o_{\mathbb{P}}(1)
\)。
\par
当 \(b<b_{\mathrm{crit}}(\beta)\) 时，取任意 \(\varepsilon\in\bigl(0,(1-\beta)H(w)-b\log 2\bigr)\)，并令事件
\[
\mathcal{E}_N:=\left\{\left|\frac{1}{N}\log\mathcal{M}_k-(1-\beta)H(w)\right|\le \varepsilon\right\}.
\]
由 \(\frac{1}{N}\log\mathcal{M}_k=(1-\beta)H(w)+o_{\mathbb{P}}(1)\) 得 \(\mathbb{P}(\mathcal{E}_N)\to 1\)。
在 \(\mathcal{E}_N\) 上对充分大 \(N\) 有 \(2^B\le \mathcal{M}_k\)，从而 \(\min(1,2^B/\mathcal{M}_k)=2^B/\mathcal{M}_k\)。
再由 \(B=\lfloor bN\rfloor\) 得 \(2^B=\exp(bN\log 2+O(1))\)，且在 \(\mathcal{E}_N\) 上有
\[
\exp\!\bigl(-N((1-\beta)H(w)+\varepsilon)\bigr)
\le
\mathcal{M}_k^{-1}
\le
\exp\!\bigl(-N((1-\beta)H(w)-\varepsilon)\bigr).
\]
对 \(\mathrm{Succ}^{\mathrm{opt}}_{\beta,b}(N)=\EE[\min(1,2^B/\mathcal{M}_k)]\) 取期望并用 \(\mathbb{P}(\mathcal{E}_N^c)=o(1)\)，得到
\[
\exp\!\bigl(N(b\log 2-(1-\beta)H(w)-\varepsilon)+o(N)\bigr)
\le
\mathrm{Succ}^{\mathrm{opt}}_{\beta,b}(N)
\le
\exp\!\bigl(N(b\log 2-(1-\beta)H(w)+\varepsilon)+o(N)\bigr).
\]
令 \(\varepsilon\downarrow 0\) 即得所示负指数率。
\par
当 \(b\ge b_{\mathrm{crit}}(\beta)\) 时，取任意 \(\varepsilon>0\) 并令
\[
\mathcal{E}_N(\varepsilon):=\left\{\left|\frac{1}{N}\log\mathcal{M}_k-(1-\beta)H(w)\right|\le \varepsilon\right\},
\]
则 \(\mathbb{P}(\mathcal{E}_N(\varepsilon))\to 1\)。
由 \(B=\lfloor bN\rfloor\) 与 \(b\log 2\ge (1-\beta)H(w)\) 可知在 \(\mathcal{E}_N(\varepsilon)\) 上有
\[
\min\!\left(1,\frac{2^B}{\mathcal{M}_k}\right)
\ge
\exp(-\varepsilon N+O(1)).
\]
取期望并用 \(\mathbb{P}(\mathcal{E}_N(\varepsilon))\to 1\) 得
\[
\mathrm{Succ}^{\mathrm{opt}}_{\beta,b}(N)\ge \exp(-\varepsilon N+o(N)),
\]
从而 \(\liminf_{N\to\infty}\frac{1}{N}\log \mathrm{Succ}^{\mathrm{opt}}_{\beta,b}(N)\ge -\varepsilon\)。
令 \(\varepsilon\downarrow 0\) 得 \(\liminf\ge 0\)；再由 \(\mathrm{Succ}^{\mathrm{opt}}_{\beta,b}(N)\le 1\) 得 \(\limsup\le 0\)，故指数率为 \(0\)。
证毕。
\end{proof}

\begin{conclusion}[超临界侧的失败指数：\texorpdfstring{$b>b_{\mathrm{crit}}(\beta)$}{b>bcrit(beta)} 时失败概率的精确变分式]\label{con:xi-terminal-supercritical-failure-exponent-variational}
沿用结论 \ref{con:xi-terminal-linear-strong-converse-query-sideinfo} 的设定。
定义失败概率
\[
\mathrm{Fail}^{\mathrm{opt}}_{\beta,b}(N):=1-\mathrm{Succ}^{\mathrm{opt}}_{\beta,b}(N).
\]
设 \(n\) 固定且 \(w\) 固定，并取 \(b>b_{\mathrm{crit}}(\beta)\)。
令
\[
I_\beta(r):=
\begin{cases}
(1-\beta)D_{\mathrm{KL}}(r\|w)+\beta D_{\mathrm{KL}}\!\left(\dfrac{w-(1-\beta)r}{\beta}\,\middle\|\,w\right),
& \text{若 }\ \dfrac{w-(1-\beta)r}{\beta}\in\Delta(X),\\[10pt]
+\infty, & \text{否则},
\end{cases}
\]
则极限
\[
E_{\mathrm{fail}}(\beta,b)
:=
\lim_{N\to\infty}-\frac{1}{N}\log \mathrm{Fail}^{\mathrm{opt}}_{\beta,b}(N)
\]
存在，并满足精确变分式
\[
\boxed{\;
E_{\mathrm{fail}}(\beta,b)
=
\inf\left\{I_\beta(r):\ (1-\beta)H(r)\ge b\log 2\right\}.\;}
\]
特别地，\(E_{\mathrm{fail}}(\beta,b)>0\) 且随 \(b\) 严格增大。
\end{conclusion}
\begin{proof}[证明]
由结论 \ref{con:xi-terminal-linear-strong-converse-query-sideinfo} 的表示，
\[
\mathrm{Fail}^{\mathrm{opt}}_{\beta,b}(N)
=
\mathbb{E}\!\left[\Bigl(1-\frac{2^B}{\mathcal{M}_k}\Bigr)_+\right].
\]
同样由 Stirling 展开，
\(
\frac{1}{N}\log\mathcal{M}_k=(1-\beta)H(R_N)+o_{\mathbb{P}}(1)
\)。
当事件 \((1-\beta)H(R_N)\ge b\log 2+\varepsilon\) 发生时，有
\(
2^B/\mathcal{M}_k\le \exp(-\varepsilon N+o(N))
\)，从而条件失败趋于 \(1\)；
当 \((1-\beta)H(R_N)\le b\log 2-\varepsilon\) 时，条件失败为 \(0\)（对充分大 \(N\)）。
因此
\[
\mathrm{Fail}^{\mathrm{opt}}_{\beta,b}(N)
=
\mathbb{P}\bigl((1-\beta)H(R_N)\ge b\log 2\bigr)\cdot(1+o(1)).
\]
而 \((R_N)_{N}\) 为超几何剩余经验分布，满足大偏差原理；其速率函数正为上式所定义的 \(I_\beta\)。
应用收缩原理即得
\[
\lim_{N\to\infty}-\frac{1}{N}\log \mathrm{Fail}^{\mathrm{opt}}_{\beta,b}(N)
=\inf\{I_\beta(r):(1-\beta)H(r)\ge b\log 2\}.
\]
严格单调性由约束集合随 \(b\) 严格收缩及 \(I_\beta\) 的严格凸性给出。证毕。
\end{proof}

\begin{conclusion}[最小化器的 escort--KKT 结构：失败指数的一参幂倾斜刻画与灵敏度公式]\label{con:xi-terminal-failure-exponent-escort-kkt}
在结论 \ref{con:xi-terminal-supercritical-failure-exponent-variational} 的条件下，进一步假设 \(w\) 全支撑且非均匀，并取 \(b>b_{\mathrm{crit}}(\beta)\)。
则：
\begin{enumerate}
  \item 变分问题
  \(
  \inf\{I_\beta(r):(1-\beta)H(r)\ge b\log 2\}
  \)
  存在唯一最小化器 \(r^\star\in\Delta(X)\)，且必位于边界 \((1-\beta)H(r^\star)=b\log 2\)。
  \item 记
  \[
  q^\star:=\frac{w-(1-\beta)r^\star}{\beta}\in\Delta(X),
  \]
  则存在唯一拉格朗日乘子 \(\theta^\star>0\) 使得 \(q^\star\) 与 \(r^\star\) 满足幂倾斜（escort）关系
  \[
  \boxed{\;
  q^\star(x)\ \propto\ r^\star(x)^{\,1+\theta^\star},\qquad x\in X.\;}
  \]
  \item 失败指数满足
  \[
  E_{\mathrm{fail}}(\beta,b)=I_\beta(r^\star)
  =(1-\beta)D_{\mathrm{KL}}(r^\star\|w)+\beta D_{\mathrm{KL}}(q^\star\|w),
  \]
  且 \(b\mapsto E_{\mathrm{fail}}(\beta,b)\) 在超临界侧可微，其导数由乘子给出：
  \[
  \boxed{\;\frac{\partial}{\partial b}E_{\mathrm{fail}}(\beta,b)=\theta^\star\log 2.\;}
  \]
\end{enumerate}
\end{conclusion}
\begin{proof}[证明]
可行域 \(\{r:(1-\beta)H(r)\ge b\log 2\}\) 为凸集（因 \(H\) 凹，超水平集凸），而 \(I_\beta\) 在全支撑内为严格凸函数，故最小化器唯一且必满足约束取等（否则可沿指向 \(w\) 的线段降低 \(I_\beta\) 并保持可行）。
\par
对带不等式约束的凸优化应用 KKT 条件：存在唯一 \(\theta^\star\ge 0\) 使得
\[
\nabla I_\beta(r^\star)+\theta^\star\,\nabla\bigl(b\log 2-(1-\beta)H(r^\star)\bigr)=\lambda\,\mathbf 1
\]
（\(\lambda\) 为归一化约束的乘子）。
在内点处 \(\nabla H(r)=-\log r-\mathbf 1\)，并且由 \(q^\star=(w-(1-\beta)r^\star)/\beta\) 可得
\(
\nabla I_\beta(r^\star)=(1-\beta)(\log r^\star-\log q^\star)
\)
（逐坐标计算）。
代入并吸收常数项，得到
\[
(\,1+\theta^\star)\log r^\star-\log q^\star=\mathrm{const},
\]
即 \(q^\star(x)\propto r^\star(x)^{1+\theta^\star}\)。
\par
灵敏度公式为拉格朗日乘子敏感性结论的直接应用：将约束写为 \(g_b(r):=b\log 2-(1-\beta)H(r)\le 0\)，其对参数 \(b\) 的导数为 \(\partial_b g_b=\log 2\)，故 \(\partial_b E_{\mathrm{fail}}(\beta,b)=\theta^\star\log 2\)。
证毕。
\end{proof}

\begin{conclusion}[Fibonacci 立方超平面统计的闭式：端点 \texorpdfstring{$\Theta$}{Theta}-类为唯一最大，给出体相--边界的几何判别]\label{con:xi-terminal-fibcube-hyperplane-theta-stats}
将 \(\Gamma_n\) 视为 partial cube，设 \(\Theta_i\) 为第 \(i\) 坐标翻转的 \(\Theta\)-等价类（超平面）。
则对任意 \(1\le i\le n\) 有严格闭式
\[
|\Theta_i|=F_i\,F_{n-i+1}.
\]
特别地，端点类 \(|\Theta_1|=|\Theta_n|=F_n\) 为最大 \(\Theta\)-类，并且当 \(1<i<n\) 时严格有 \(|\Theta_i|<F_n\)。
\end{conclusion}
\begin{proof}[证明]
见定理 \ref{thm:pom-fibcube-theta-class-size}。
\end{proof}
